% Laget automatisk av elevark/lag_elevark.py fra sim/tidslinje.html. Endringer her blir overskrevet.
\documentclass[11pt]{article}
\usepackage{elevark}
\begin{document}
\elevarkhode{Fysikk 1 og 2 · Fysikkens historie}{Fysikkens historie}{https://fysikklab.no/sim/tidslinje.html}{Simuleringen}

Fra Arkimedes i badekaret til gravitasjonsbølger fra svarte hull. Dra i tidslinjen, zoom inn der det skjer mye, og klikk på en hendelse for å lese om den, se en liten animasjon og finne simuleringer og forsøk som hører til. Noen hendelser er pensum i Fysikk 1 eller Fysikk 2, andre går utenfor.

\section*{Oppgaver}

\begin{oppgaveliste}
\oppgave \fase{gjett}{Gjett}Før du ser på tidslinjen: Hvilken kom først, Newtons gravitasjonslov eller oppdagelsen av elektronet? Hvor mange år tror du det gikk mellom dem?
\svar
\oppgave \fase{test}{Undersøk}Zoom inn på årene 1895–1935. Hvor mange hendelser finner du? Hvorfor tror du det skjedde så mye i fysikken akkurat da?
\svar
\oppgave \fase{test}{Undersøk}Finn de to hendelsene fra 1905. Hva har de til felles, og hvem stod bak?
\svar
\oppgave \fase{forklar}{Forklar}Finn to hendelser der en måling avgjorde mellom to teorier. Hvilke teorier sto mot hverandre, og hva viste målingen?
\svar
\oppgave \fase{forklar}{Forklar}Velg «Utenfor pensum». Velg én hendelse og forklar med egne ord hvorfor den var viktig for fysikken.
\svar
\oppgave \fase{test}{Undersøk}Finn tre hendelser der nordiske forskere var med. Hva oppdaget de?
\svar
\oppgave \fase{utfordring}{Utfordring}Mange av de nyeste hendelsene ble gjort av tusenvis av forskere i store samarbeid. Finn eksempler på tidslinjen, og diskuter hvorfor moderne fysikk ofte krever internasjonalt samarbeid.
\svar
\end{oppgaveliste}

\end{document}
